% Clase 27 — voltaje del generador y corriente en régimen permanente:
% V = V0 cos(omega t), I = (V0/|Z|) cos(omega t - theta). Dibujado con
% theta = pi/4 y |Z| tal que la amplitud de I sea 0.6 V0 (unidades arbitrarias):
% la corriente está atrasada theta respecto del voltaje.
\begin{tikzpicture}
  \begin{axis}[
    emfig, width=11.5cm, height=5.6cm,
    xmin=0, xmax=12.566, ymin=-1.3, ymax=1.55,
    xtick={0,3.1416,6.2832,9.4248,12.566},
    xticklabels={$0$,$\pi$,$2\pi$,$3\pi$,$4\pi$},
    ytick={-1,0,1},
    xlabel={$\omega t$},
    legend style={at={(0.5,-0.4)}, anchor=north, legend columns=2,
      /tikz/every even column/.append style={column sep=1em}},
  ]
    \addplot[figblue, domain=0:12.566, samples=200] {cos(deg(x))};
    \addlegendentry{$V(t)/V_0$}
    \addplot[figred, domain=0:12.566, samples=200] {0.6*cos(deg(x-0.7854))};
    \addlegendentry{$I(t)$, escala arbitraria}
    \addplot[figaux, forget plot] coordinates {(6.2832,-1.3) (6.2832,1.3)};
    \addplot[figaux, forget plot] coordinates {(7.0686,-1.3) (7.0686,1.3)};
    \draw[figgray, line width=0.6pt, {Latex[length=1.5mm]}-{Latex[length=1.5mm]}]
      (axis cs:6.2832,1.15) -- (axis cs:7.0686,1.15);
    \node[font=\small, anchor=south] at (axis cs:6.676,1.15) {$\theta$};
  \end{axis}
\end{tikzpicture}
